El análisis se construyó a partir de una separación astrofísica de las variables. En lugar de usar todas las columnas disponibles del catálogo al mismo tiempo, las variables se organizaron en espacios con significado físico: propiedades corporales del planeta, propiedades orbitales, propiedades térmicas y combinaciones entre ellas. Esta decisión es importante porque Mapper es sensible a la escala, a la redundancia entre variables y a los valores faltantes. Por ello, cada espacio de variables responde una pregunta distinta sobre la población de exoplanetas.

Las variables físicas describen qué tipo de objeto planetario se está observando. En este bloque se usan principalmente \texttt{pl\_rade}, \texttt{pl\_bmasse} y, en algunos espacios, \texttt{pl\_dens}. El radio permite distinguir planetas pequeños, sub-Neptunos, Neptunos y gigantes; la masa aporta información sobre escala gravitacional y composición posible; y la densidad resume la relación entre masa y volumen. En conjunto, estas variables aproximan la naturaleza física del planeta: si es compatible con un cuerpo rocoso, un planeta con envoltura gaseosa o un gigante gaseoso.

Las variables orbitales describen el entorno dinámico del planeta, es decir, dónde vive dentro de su sistema. En los artefactos reconciliados de este proyecto, el espacio orbital usa \texttt{pl\_orbper} y \texttt{pl\_orbsmax}. El periodo orbital indica cuánto tarda el planeta en completar una órbita, mientras que el semieje mayor aproxima su distancia media a la estrella. Estas variables permiten separar, por ejemplo, planetas de periodo corto y órbitas cercanas de planetas con órbitas más amplias o periodos largos.

Las variables térmicas describen el ambiente energético del planeta. El espacio \texttt{thermal} usa \texttt{pl\_insol} y \texttt{pl\_eqt}. La insolación mide cuánta radiación recibe el planeta en comparación con la Tierra, mientras que la temperatura de equilibrio estima la temperatura esperada bajo un modelo simplificado de balance radiativo. Estas variables son útiles para distinguir planetas muy calientes, calientes, templados o fríos. Sin embargo, en este proyecto el bloque térmico debe interpretarse con cautela, porque mostró una dependencia mucho mayor de imputación que el bloque orbital.

Los espacios de variables existentes en los artefactos reconciliados son:
\[
\texttt{joint},\quad
\texttt{joint\_no\_density},\quad
\texttt{orb\_thermal},\quad
\texttt{orbital},\quad
\texttt{phys\_density},\quad
\texttt{phys\_min},\quad
\texttt{thermal}.
\]
Estos espacios permiten comparar distintas maneras de mirar el mismo catálogo. \texttt{phys\_min} usa las variables físicas principales sin densidad; \texttt{phys\_density} agrega densidad como control físico; \texttt{orbital} se concentra en la escala orbital; \texttt{thermal} resume el ambiente energético; \texttt{orb\_thermal} combina órbita y energía recibida; y los espacios \texttt{joint} y \texttt{joint\_no\_density} permiten evaluar la estructura conjunta con y sin densidad.

\begin{table}[H]
\centering
\small
\begin{tabular}{p{0.25\textwidth}p{0.35\textwidth}p{0.30\textwidth}}
\toprule
\textbf{Espacio} & \textbf{Variables} & \textbf{Interpretación principal} \\
\midrule
\texttt{phys\_min} & \texttt{pl\_rade}, \texttt{pl\_bmasse} & Tamaño y masa del planeta \\
\texttt{phys\_density} & \texttt{pl\_rade}, \texttt{pl\_bmasse}, \texttt{pl\_dens} & Propiedades físicas incluyendo densidad \\
\texttt{orbital} & \texttt{pl\_orbper}, \texttt{pl\_orbsmax} & Escala orbital: periodo y distancia media \\
\texttt{thermal} & \texttt{pl\_insol}, \texttt{pl\_eqt} & Ambiente energético del planeta \\
\texttt{orb\_thermal} & \texttt{pl\_orbper}, \texttt{pl\_orbsmax}, \texttt{pl\_insol}, \texttt{pl\_eqt} & Relación entre órbita y ambiente térmico \\
\texttt{joint\_no\_density} & Físicas, orbitales y térmicas sin \texttt{pl\_dens} & Estructura conjunta sin densidad derivada \\
\texttt{joint} & Físicas, orbitales y térmicas con \texttt{pl\_dens} & Estructura conjunta incluyendo densidad \\
\bottomrule
\end{tabular}
\caption{Espacios de variables usados para construir los grafos Mapper reconciliados.}
\label{tab:feature_spaces_context}
\end{table}

Un punto importante es que el espacio orbital seleccionado no incluye excentricidad. El grafo \texttt{orbital\_pca2\_cubes10\_overlap0p35} fue construido con las variables registradas en su JSON de configuración: \texttt{pl\_orbper} y \texttt{pl\_orbsmax}. Por lo tanto, en los artefactos reconciliados, \texttt{pl\_orbeccen} no forma parte de las coordenadas del Mapper orbital actual. Esto no significa que la excentricidad sea irrelevante; al contrario, físicamente describe la forma de la órbita y puede distinguir órbitas circulares de órbitas altamente elípticas. La decisión aquí es metodológica: como el objetivo del reporte es auditar e interpretar los grafos ya generados, incluir excentricidad requeriría construir un nuevo espacio orbital y correr nuevos Mappers. Por ello, \texttt{pl\_orbeccen} se considera una extensión futura de sensibilidad orbital, no una variable activa de la interpretación principal.

También es necesario tratar la densidad con cuidado. La densidad planetaria es una variable físicamente útil porque relaciona masa y radio, y puede ayudar a distinguir cuerpos rocosos de planetas gaseosos o de baja densidad. Sin embargo, cuando \texttt{pl\_dens} se usa junto con \texttt{pl\_bmasse} y \texttt{pl\_rade}, no debe interpretarse como evidencia completamente independiente. En la auditoría de derivaciones, 6199 valores de \texttt{pl\_dens} fueron derivados a partir de masa y radio; antes de la derivación había 6273 valores faltantes y después quedaron 74. Esta dependencia algebraica justifica comparar explícitamente espacios con y sin densidad. Por eso \texttt{phys\_density} y \texttt{joint} funcionan como controles de sensibilidad, mientras que \texttt{phys\_min} y \texttt{joint\_no\_density} permiten observar la estructura sin incorporar una variable derivada.

Además de las variables usadas como coordenadas de Mapper, el proyecto conserva variables de auditoría. \texttt{discoverymethod}, \texttt{disc\_year} y \texttt{disc\_facility} no describen al planeta como objeto físico, sino el proceso mediante el cual fue descubierto. Por esa razón no se usan como variables de entrada del Mapper. Incluirlas directamente en la geometría podría producir grafos que agrupen planetas por su historia observacional en vez de por sus propiedades físicas u orbitales. En cambio, se usan después como metadata externa para preguntar si las regiones del grafo están enriquecidas por método, época o instalación de descubrimiento.

Esta distinción es clave porque los métodos de detección no observan de manera uniforme toda la población planetaria. El método de tránsito favorece planetas con periodos cortos y radios grandes, porque producen caídas de brillo más frecuentes y detectables. La velocidad radial favorece planetas masivos, especialmente si inducen variaciones medibles en la estrella. La imagen directa favorece objetos jóvenes, masivos, luminosos y separados de su estrella. Por lo tanto, una región Mapper dominada por un solo método de descubrimiento puede reflejar tanto una estructura física real como una huella del proceso de observación.

Finalmente, las etiquetas interpretativas usadas en el análisis, como \texttt{radius\_class}, \texttt{density\_class}, \texttt{orbit\_class}, \texttt{thermal\_class} o \texttt{candidate\_population}, deben leerse como apoyos heurísticos y no como una taxonomía definitiva. Su función es ayudar a describir qué tipo de planetas aparecen en cada nodo o componente del grafo. Una región será más convincente si combina coherencia física u orbital, baja imputación y baja asociación con \texttt{discoverymethod}. Una región será más sospechosa si está dominada por un método de descubrimiento, una facilidad específica o una alta proporción de valores imputados.